\chapter{Laboratory workflows and reproducibility}
\label{dnad26:chapter}
\ChapterGoals{%
\item turn an abstract claim about a molecular computation into a testable record;
\item distinguish design identity, physical identity and measured identity of a molecular sample;
\item bind evidence to exact file contents and derive an append-only invariant for a ledger;
\item propagate preparation uncertainty and choose controls that test mechanisms;
\item distinguish technical repeats from independent replicates, and preserve raw observations.}
{Chapters~10 and~25; hashes are introduced as needed.}
A diagram specifies what a molecular system is intended to do. A laboratory record
must establish which physical realization was actually tested, what was observed,
and how those observations became a conclusion. These are different questions.
Chapter 25 recovered a digital file under declared coding assumptions; it did not
establish that an ordered strand was intact, an instrument was calibrated, or an
independent preparation would behave the same way.

Our objective is to build a small, executable \Term{evidence ledger}. We derive its
transition rules, implement them, try to break them, and examine what they cannot
verify. The ledger is a provenance teaching artifact, not an operating protocol,
a certified electronic laboratory notebook, or evidence that a wet-lab experiment
was performed. All example artifact bytes and replicate outcomes are assigned
synthetic teaching data.

\section{From an abstract claim to a testable record}
Before collecting evidence, write a claim with an intervention, outcome,
denominator, comparison and experimental unit. For example, a design revision
improves the fraction of correctly classified assemblies across independently
prepared reactions. This specifies more than saying that the design works.
A fraction of correct objects, a fraction of functional reaction vessels and
a fraction of recoverable files are different estimands.

A design freeze should bind the primary outcome and its classifier, control roles,
planned independent units, exclusion rules, stopping or sampling plan, and the
boundary between confirmatory and exploratory analysis. Freezing does not prohibit
learning. An exploratory revision receives a new identity and remains valuable;
it must not silently replace the prospectively specified question.

For molecular implementations, a sequence manifest needs identities, roles and
orientation, with an explicit convention for complement and reverse complement.
Chapter 24's resource accounting and Chapter 25's sequence coding do not supply
a universal laboratory recipe. Buffer, instrument and material requirements depend
on the mechanism and validated local procedures. The DNA-origami primer provides
context for design, synthesis and characterization \cite{Dey2021}; this chapter
does not claim to reproduce its laboratory methods.

\section{Six boundaries, not six decorative labels}
Let $q_i$ be one of six ordered stages:
design frozen, materials verified, run recorded, observations frozen, analysis
frozen, and released. A history is $H_k=(E_0,\ldots,E_{k-1})$, not an isolated label
saying released. Define the required new artifact roles
\[
\begin{aligned}
R_0&=\{\mathrm{design,sequences,protocol}\},\\
R_1&=\{\mathrm{materials}\},\\
R_2&=\{\mathrm{execution,controls,randomization}\},\\
R_3&=\{\mathrm{raw,calibration,blinding}\},\\
R_4&=\{\mathrm{analysis\_spec,analysis\_code,derived}\},\\
R_5&=\{\mathrm{release\_metadata}\}.
\end{aligned}
\]
A stage-$i$ event must retain all roles in $\bigcup_{j=0}^{i}R_j$. Presence is a
necessary structural prerequisite, not a scientific judgment about an artifact's
contents. A file labeled calibration can still contain a failed calibration.
\Plate{01-state}{Each audited boundary adds evidence while retaining earlier
identities and artifacts. Cumulative prerequisites prevent a release label from
standing in for the complete history.}{dnad26:state}

The demo binds three artifacts at design freeze, then four, seven, ten,
thirteen and fourteen. These counts follow from this deliberately small
schema, not from a universal reporting standard.
\begin{center}\input{\Generated/stages.tex}\end{center}
The release gate has two distinct parts: machine validation of the ledger, and
human/domain validation of the experiment and inference. Neither should be
reported as the other.

\section{Identity is a hierarchy of physical commitments}
A \Term{design identity} names an abstract mechanism. A sequence-set identity names
its specific strand realization. A material identity names the relevant acquired
or produced lot and its provenance. A preparation identity names a specific
physical combination, including its relation to an independent experimental unit.
\Plate{02-identity}{A design can have multiple sequence realizations; a sequence
realization can have different material lots. Multiple measurements of one
preparation are not new preparations.}{dnad26:identity}

Consider three changes. Correcting a base in the sequence manifest changes the
sequence-set identity even if the conceptual gate is unchanged. Reordering the
same sequences from a new lot changes material identity without necessarily
changing the sequence set. Imaging a second field from the same prepared vessel
changes the observation identity, not the preparation identity.

Our frozen \texttt{Plan} binds run, design, sequence set, material, preparation,
protocol version, independent-unit identity, distinct control identities and a
blinding policy. Identifiers are nonempty stripped strings. Control identities
must be distinct; the two-control minimum encodes this teaching contract, not
a sufficient experimental design for every assay. A material substitution or
sequence revision after the initial freeze requires a new linked run. A correction
is recorded with its reason and provenance; previous evidence is not overwritten.

\section{Bind actual bytes, not just plausible-looking hashes}
For artifact bytes $B$ define
\[
A=(\mathrm{role,id,SHA256}(B),|B|).
\]
Validation retrieves bytes by identity and checks type, size and digest. Checking
only that a string has 64 hexadecimal characters would accept a missing file.
Checking only size would accept a same-length mutation.

The implementation uses an immutable tuple of frozen artifact records and rejects
duplicate artifact identities. Immutability of the record does not make the external
byte store immutable: that is why every validation rechecks the actual bytes.
\Listing{verify_artifacts}
The function assumes \texttt{Artifact}, \texttt{ROLES}, \texttt{identifier} and
\texttt{sha} from the complete runnable module. Each listing is the literal tested
function, not abbreviated pseudocode. The full file also defines the plan and event
dataclasses, constructors and validation helpers.

Our event digest uses a declared local JSON profile: sorted keys, compact separators,
ASCII escaping, no non-finite JSON numbers and ASCII bytes. This is a reproducibility
contract for these Python records, not conformance to a cross-language canonical-JSON
standard. Another implementation must reproduce the specified byte representation
or adopt a separately versioned serialization.

\section{Derive the append-only invariant}
Write $P(E)$ for the frozen plan, $A(E)$ for its artifact set, and $D(E)$ for its
ordered deviations. A valid adjacent pair must satisfy
\[
\begin{aligned}
q(E_i)&=q_i,&
\mathrm{previous}(E_i)&=\mathrm{digest}(E_{i-1}),\\
P(E_i)&=P(E_{i-1}),&
A(E_{i-1})&\subseteq A(E_i),\\
D(E_{i-1})&\preceq D(E_i).
\end{aligned}
\]
Here set inclusion compares the entire artifact record, not just its identity;
$\preceq$ means ordered-prefix inclusion. The genesis predecessor is a fixed
all-zero digest. The event's own digest must match its serialized content.
\Listing{validate_history}

This yields a short inductive argument. The genesis event fixes the plan and its
first artifacts. Each validated successor preserves them and appends evidence.
Therefore every valid prefix retains all earlier bound records and deviations.
The statement is conditional on the supplied history and byte store. It says
nothing about whether an omitted real-world event occurred.
\Listing{advance}
The constructor advances exactly one stage, validates both the existing history
and its actual bytes, and rejects a second release. A restarted instrument can
be recorded as a deviation before observations freeze. Its disposition belongs
in the analysis specification; appending the note does not automatically excuse
or invalidate the run.

\subsection{The coherent-forgery counterexample}
Suppose an adversary changes artifact bytes, updates their digests, rewrites the
events and recomputes every predecessor link. The resulting internally consistent
history can pass this validator. A hash chain detects inconsistency relative to
a trusted earlier record; it is not an identity authority or an unforgeable timestamp.
External anchoring, authenticated signatures, access control and retention policies
address different threats and need their own design. Even a signed false measurement
remains false. Describe this system as byte binding and structural provenance,
not cryptographic proof of experimental truth.

\section{Preparation quantities: derive, convert, qualify}
For concentration $c$ in mol/L and volume $V$ in L, nominal amount is
\begin{equation}
n=cV.
\label{dnad26:amount}
\end{equation}
Thus $2\,\mathrm{nM}=2\times10^{-9}\,\mathrm{mol/L}$ and
$5\,\mu\mathrm L=5\times10^{-6}\,\mathrm L$ give $n=10^{-14}$ mol.
The unit conversion is part of the computation, not typographical decoration.
\Plate{03-ledger}{Nominal amount follows from compatible units. Uncertainty and
functional fraction are separate questions; the example is assigned,
not a measured preparation.}{dnad26:ledger}

If the intact functional fraction is $f$, functional amount would be $fn$ under
that model. The ledger does not measure $f$. Concentration inferred from absorbance,
supplier nominal concentration, folded fraction and incorporated fraction are not
interchangeable quantities. Record the measurement method and what it estimates.

Let concentration and volume uncertainties be $\sigma_c,\sigma_V$ with correlation
$\rho$. Linearize $n(c+\delta c,V+\delta V)$:
$\delta n\approx V\delta c+c\delta V$. Taking variance gives
\begin{equation}
\sigma_n^2\approx V^2\sigma_c^2+c^2\sigma_V^2+
2cV\rho\sigma_c\sigma_V.
\label{dnad26:uncertainty}
\end{equation}
This first-order approximation requires a defensible uncertainty model.
It omits $\delta c\,\delta V$ and need not describe
large, skewed or systematically biased errors.

For $\sigma_c=0.2$ nM, $\sigma_V=0.1\,\mu$L and $\rho=0$, the example yields
$\sigma_n=1.0198039027\times10^{-15}$ mol, approximately $10.2\%$ relative
uncertainty. Correlated errors change the answer; independence must not be assumed
merely because concentration and volume have different names.
\Listing{amount_uncertainty}
The code evaluates the equivalent nonnegative form
$\operatorname{hypot}(x+\rho y,\sqrt{1-\rho^2}\,y)$ with $x=V\sigma_c$,
$y=c\sigma_V$. It rejects invalid inputs and non-finite results rather than releasing
an infinite uncertainty as a usable result. This numerical arrangement is not
a higher-order physical uncertainty model.

\section{Controls test mechanisms; identities alone do not}
A negative control asks what happens without the proposed intervention or required
mechanism. A positive control asks whether the measurement can recover an expected
signal under a validated condition. A process control probes a specific upstream
failure such as sample handling or contamination. A no-input gate, a known readout
standard and a sample-handling blank answer different questions.

Declare the control's prediction, denominator and disposition rule in advance.
A failed positive control can make an apparent absence of output uninterpretable;
a successful positive control does not prove the designed gate computed correctly.
High negative-control background can invalidate a threshold-based interpretation.
Our schema verifies the existence of a control record, not whether any prediction held.

Randomized processing order can break alignment between condition and temporal
drift. Blinded classification can block expectation from affecting field selection
or borderline calls. Neither fixes an irrelevant control or a defective instrument.
Record the assignment artifact and its scope: randomizing images after preparation
does not randomize preparation-level treatment.

\section{Three preparations are not three hundred replications}
Let $y_{rj}$ be observation $j$ in independently prepared unit $r$, with $m_r$
observations. For a preparation-weighted estimand define
\[
\bar y_r=\frac{1}{m_r}\sum_j y_{rj},\qquad
\hat\mu=\frac1R\sum_r\bar y_r,\qquad
\widehat{\operatorname{Var}}(\hat\mu)=
\frac{\sum_r(\bar y_r-\hat\mu)^2}{R(R-1)}.
\]
This is the sample variance of independent preparation means divided by $R$.
It requires independent units with an appropriate common-population interpretation;
it is not an exact small-sample confidence interval or a universal hierarchical model.
\Plate{04-replicates}{Treatment assigned to preparations creates shared upstream
variation. Increasing object counts does not increase the number of independently
prepared units.}{dnad26:replicates}
\Listing{cluster_mean_variance}

Assign 100 binary objects to each of three preparations, with preparation means
$(0,1,1)$. Then $\hat\mu=2/3$ and preparation-mean variance is $1/9$.
Treating all 300 objects as independent instead gives
$\hat\mu(1-\hat\mu)/299=0.0007432181$, smaller by a factor of $149.5$.
The example deliberately creates perfect within-preparation agreement to expose
the independence error. It is not an estimate of laboratory variability.

When $m_r$ differ, an object-weighted mean and a preparation-weighted mean can differ.
For means $(0,1)$ with counts $(1,999)$, equal preparation weighting yields $0.5$,
while pooling objects yields $0.999$. Neither is automatically the correct target.
Specify the estimand before selecting an uncertainty method.

\section{Preserve observation before classification}
An image, fluorescence trace or sequencing output is not yet a scientific success
label. A classifier maps a defined raw artifact and calibration to a reported
variable. Freeze image selection, segmentation, thresholds, handling of ambiguous
objects and denominator conventions. Otherwise a classifier revision changes
reported accuracy while no molecular mechanism changed.
\Plate{05-measure}{Raw observations, calibrated variables and success labels are
distinct artifacts. Versioning the mapping permits reanalysis without rewriting
the physical preparation.}{dnad26:measure}

Our blinding contract requires coded identities at observation freeze when the plan
requires blinding. Later analysis may unblind after that recorded boundary; the
history preserves the earlier blinded event. If blinding is infeasible, the plan
must include an explicit reason rather than silently setting a Boolean false.
A real study may need blinded analysis through additional stages and authenticated
custody of the code key.

Retain prespecified results even when a deviation motivates a sensitivity
analysis. An alternative exclusion or threshold is a new analytical artifact with
its own justification. It must not erase the original analysis or disguise
a post-hoc decision as a frozen rule.

\section{Release a usable evidence bundle}
\Plate{06-release}{Machine-valid provenance is only one release gate. Scientific
interpretation, access policy and external replication remain separate assessments.}
{dnad26:release}
A release should bind the question, identity hierarchy, raw evidence, calibration,
control outcomes, deviations, analysis specification, code/environment versions and
derived tables and figures. A reader should reconstruct each plotted value
and discover why an excluded observation was excluded.

FAIR means findable, accessible, interoperable and reusable under appropriate
metadata and stewardship \cite{Wilkinson2016}. It does not require unrestricted
public data: authentication and authorization may be compatible with accessibility.
Record the access mechanism, license and retention policy. A ZIP with unnamed
images and no denominator is not reusable simply because it downloads.

\subsection{Four levels of evidence}
Computational reproduction reruns the frozen analysis on the same bytes.
Analytical robustness asks how reasonable analysis choices affect the conclusion.
Experimental replication obtains new independently prepared physical evidence.
External replication tests the claim at another team or site.
These gates establish different things. Passing the software tests below establishes
none of the last three by itself.

\section{Independent tests and deliberate counterexamples}
The tests recompute event hashes through a separately constructed JSON body, not the
module's serialization helper. They enumerate illegal stage orders, remove required
roles even after recomputing a valid digest, mutate same-length bytes, substitute a
renamed artifact, change the frozen plan, delete a deviation and break a predecessor
link. This separates content-integrity checks from workflow checks.

A rehashed isolated release is rejected because it is not genesis.
A rehashed artifact replacement can pass an individual event's role check but fail
history monotonicity. The blinding test rejects an unblinded observation freeze
and allows later unblinding. Hand-derived covariance and preparation-variance
oracles test the arithmetic independently.
The suite also covers malformed digests, duplicate identities,
Boolean values masquerading as numbers, invalid quantities and overflow.

The demonstration writes fourteen assigned text artifacts, not experimental files.
Its six valid events establish execution of the stated ledger contract.
Laboratory competence, independent scientific review, assay validation and replication
remain open. This limitation is part of the result, not fine print.

\begin{KeyIdeas}
\begin{itemize}\setlength{\itemsep}{2pt}
\item A design states intent; a record must establish what was actually made,
measured and concluded, and these are separate boundaries.
\item A sequence name, an ordered oligonucleotide and a verified molecule are
different levels of identity.
\item Cryptographic hashes bind evidence to exact bytes; an append-only ledger
makes later changes visible.
\item Preparation quantities carry uncertainty, $\sigma_n^2\approx V^2\sigma_c^2+c^2\sigma_V^2+2cV\rho\sigma_c\sigma_V$, which must be reported.
\item Controls test mechanisms: a negative control can show that a signal needs the
intended input, not merely that the instrument works.
\item Three readings of one preparation are not three experiments; preserve raw
observations before classification.
\end{itemize}
\end{KeyIdeas}

\section{Twenty worked exercises}
\begin{enumerate}
\item Same conceptual gate, corrected base: which identity changes?
\textbf{Answer:} sequence-set identity; the design identity may remain.
\item New lot with unchanged sequences?
\textbf{Answer:} new material identity and new run/preparation provenance.
\item Second image from one vessel: a new independent unit?
\textbf{Answer:} no, a new observation nested within that preparation.
\item What roles are required at observation freeze?
\textbf{Answer:} all ten roles accumulated from stages zero through three.
\item Why reject an isolated correctly hashed release?
\textbf{Answer:} a valid history begins at design freeze and advances in order.
\item Can a well-formed digest establish that bytes exist?
\textbf{Answer:} no; retrieve bytes and verify length and digest.
\item A same-length byte mutation preserves size. Which check catches it?
\textbf{Answer:} recomputing SHA-256 against the bound artifact digest.
\item Why compare full artifact records rather than identity sets?
\textbf{Answer:} the same identity could otherwise name altered bytes or a new role.
\item Can an internally rebuilt chain authenticate its author?
\textbf{Answer:} no; recomputed links can describe a coherent forgery.
\item Convert 2 nM in $5\,\mu$L to mol.
\textbf{Answer:} $(2\times10^{-9})(5\times10^{-6})=10^{-14}$ mol.
\item Derive the covariance term in Eq.~\ref{dnad26:uncertainty}.
\textbf{Answer:} variance of $V\delta c+c\delta V$ contains
$2cV\operatorname{Cov}(\delta c,\delta V)=2cV\rho\sigma_c\sigma_V$.
\item Compute $\sigma_n$ for the stated independent errors.
\textbf{Answer:} $\sqrt{(10^{-15})^2+(2\times10^{-16})^2}
=1.0198039\times10^{-15}$ mol.
\item Does $nN_A$ establish functional strand count?
\textbf{Answer:} no; intactness, availability and incorporation need separate evidence.
\item Derive variance of the equally weighted preparation mean.
\textbf{Answer:} independent preparation sample variance divided by $R$ gives
$\sum_r(\bar y_r-\hat\mu)^2/[R(R-1)]$.
\item Evaluate it for preparation means $(0,1,1)$.
\textbf{Answer:} $\hat\mu=2/3$, squared-deviation sum $2/3$, variance $1/9$.
\item Compare pooling and equal weighting for counts $(1,999)$, means $(0,1)$.
\textbf{Answer:} pooled $0.999$, equally weighted $0.5$; different targets.
\item Can a failed positive control support a claim of no gate output?
\textbf{Answer:} not without resolving whether the assay could observe output.
\item When may the example unblind?
\textbf{Answer:} after the required blinded observation-freeze boundary;
stricter study policies may delay this.
\item Must FAIR data always be unrestricted?
\textbf{Answer:} no; documented authorized access can be appropriate.
\item Design a release review.
\textbf{Rubric:} validate bytes/history, inspect controls/denominators,
reproduce derived results, assess deviations/uncertainty, document access,
and report replication and independent review separately.
\end{enumerate}

\section{Bridge to symbolic molecular simulation}
The next chapter represents molecular species and reactions as symbolic
states and rules. A simulator also needs frozen inputs, versioned semantics,
independent checks and an output trace. Reproducible symbolic execution is evidence
about a model, not by itself evidence that its molecular system operates in a vessel.
